\chapter{Fidelidad diferencial y experimentación controlada}
\label{ch:differential_experiments}
\label{ch:experimental_design}

Este capítulo define los protocolos que producirán la evidencia dinámica. No se presentan como resultados mediciones aún no realizadas. Su función es convertir H1--H3 en intervenciones donde instrumentación, movimiento y referencia puedan separarse.

\section{E0: caracterización del piso instrumental}

Con el estado físico fijo se observa
\begin{equation}
\widetilde H(t)\mid q=q_0.
\end{equation}
Se realizarán bloques dentro de una sesión, entre sesiones, después de reinicios y después de reconexiones controladas. Cuando el hardware lo permita se registrarán temperatura, ganancia, CFO, estado de reloj y telemetría. El control de escena vacía se repetirá al inicio y al final; una perturbación simulada en procesamiento verificará que la cadena de análisis puede detectarla.

La salida de E0 será una distribución de deriva y discontinuidades para cada observable. Se compararán canal complejo, magnitud, $U(1)$ por arreglo, producto espectral seleccionado y rama referenciada. Una representación se considerará más repetible sólo si el cambio excede la incertidumbre entre sesiones y no proviene de ajustar con los datos evaluados.

\section{Referencia coherente y límites de cancelación}

Para dos cadenas simultáneas con perturbación común $c(t,k)$,
$Y_s=cg_sH_s+n_s$ y $Y_r=cg_rH_r+n_r$. Un cociente regularizado es
\begin{equation}
\widetilde H_s=H_{r,\rm cal}\frac{Y_sY_r^*}{|Y_r|^2+\epsilon_r}.
\label{eq:cabled_reference_ratio}
\end{equation}
La cancelación es válida si la perturbación es común, la referencia no se anula y la relación de cadenas está calibrada. La regularización introduce sesgo y una referencia física puede contener temperatura, movimiento de cable o multitrayectoria propia. Compartir \SI{10}{MHz} o PPS facilita frecuencia/tiempo, pero no garantiza fase relativa después de un reinicio. Estas condiciones se comprobarán y no se asumirán por la hoja de datos.

\section{E1: perturbación mecánica conocida}

Un reflector se desplazará alrededor de un estado basal con posición y orientación medidas. Para cada paso $h$ se estimará
\begin{equation}
\widehat J_q^{(r)}=\frac{H(q+h\vect e_r)-H(q-h\vect e_r)}{2h},
\label{eq:finite_difference_jacobian}
\end{equation}
y su equivalente después de $V$. Se ensayarán varios $h$ para identificar un régimen donde error de truncamiento y amplificación de ruido sean controlables. La secuencia alternará signos y repetirá estados para medir histéresis y deriva.

El gemelo se calibra sólo en el basal. Sus predicciones principales son
\begin{equation}
\Delta H_{\rm sim},\qquad J_{\rm sim},
\end{equation}
comparadas con $\Delta H_{\rm real}$ y $J_{\rm real}$. Las métricas incluyen error complejo absoluto, correlación/dirección del Jacobiano, norma, valores singulares cuando $q$ es multidimensional y desempeño de una tarea de desplazamiento. Un denominador pequeño se reportará como baja sensibilidad y no como error relativo favorable.

\section{Comparación de representaciones en E1}

La comparación primaria incluye:
\begin{enumerate}
\item CSI/CFR complejo con referencia basal;
\item magnitud;
\item distancia $U(1)$ por arreglo seleccionada en S4.7-R;
\item producto espectral auxiliar R3;
\item rama coherentemente referenciada, si E0 demuestra estabilidad;
\item fusiones definidas con entrenamiento/validación.
\end{enumerate}
Para cada alternativa se medirán repetibilidad basal y respuesta física. La discriminabilidad espacial estática de S4.7-R no se reutilizará como sustituto de sensibilidad. Se incluirán geometrías previstas por la \cref{eq:bistatic_sensitivity} como favorables y casi ciegas para probar el valor predictivo del gemelo.

\section{E2: fantoma respiratorio}

El fantoma ejecutará movimientos de amplitud, frecuencia, orientación y forma de
onda conocidas. No se limitará a una sinusoide, porque ésta permite estimar una
frecuencia aun cuando el sistema distorsione la morfología. El conjunto de excitaciones
incluirá: ciclo sinusoidal; inspiración rápida con espiración lenta; el caso inverso;
amplitud modulada; pausas; dos frecuencias próximas en bloques separados; y movimiento
grueso superpuesto. Cada perfil se almacena como referencia continua $x_{\rm ref}(t)$,
junto con incertidumbre de desplazamiento y tiempo.

El estimando primario será la forma de onda $\widehat x(t)$. Para una alineación
temporal estimada exclusivamente con validación, se reportarán
\begin{align}
\operatorname{NRMSE}
&=\frac{\|\widehat{\vect x}-\vect x_{\rm ref}\|_2}
{\|\vect x_{\rm ref}-\bar x_{\rm ref}\vect 1\|_2},\\
\rho_{\max}
&=\max_{\ell\in\mathcal L_{\rm val}}
\operatorname{corr}\{\widehat x(t-\ell),x_{\rm ref}(t)\}.
\end{align}
El retardo seleccionado, el error de amplitud por ciclo, las duraciones de
inspiración y espiración, la detección de pausas, la coherencia espectral y una
distancia de deformación temporal se informarán por separado. La frecuencia
respiratoria se derivará sólo en intervalos donde una frecuencia instantánea resulte
significativa; no será sustituto de la forma de onda.

\subsection{Contrastes de fase y selección de enlaces}

Sobre la misma adquisición se compararán: magnitud; fase cruda; fase con CPO
retirada; retardo predicho más fase diferencial; retardo oráculo como techo de
alineabilidad; productos entre antenas; productos entre frecuencias; y canal
complejo coherentemente referenciado. El caso oráculo se etiqueta como límite
superior y no participa en la comparación de despliegue. Cada alternativa utilizará
idénticas ventanas y particiones.

La selección mediante SNR de comunicación se contrastará con SSNR/SCR y con
información equivalente \(F_{x\mid\zeta}\). El análisis probará tres afirmaciones
operacionales: si el retardo efectivo mejora la combinación coherente y la forma
de onda; si eliminar fase común reduce el piso sin proyectar \(J_x\); y si la
fidelidad del Jacobiano predice el error respiratorio mejor que potencia o fidelidad
estática. Estas afirmaciones son contrastes pendientes, no conclusiones del banco
dc01.

Una forma de onda mecánica recuperada no valida respiración humana ni apnea clínica. E2 determina si la cadena observa un micromovimiento conocido y si el gemelo predice condiciones favorables. La transición a E3 añadirá deformación no rígida, postura, ropa y movimiento espontáneo.

\section{Preflight simulado y criterio de paso}

Antes de la adquisición se podrán simular perturbaciones conocidas en Sionna RT, WaveVerse u otra plataforma para detectar representaciones que anulan el movimiento por construcción \citep{Hoydis2023SionnaRT,Zheng2026WaveVerse}. Este preflight es un diagnóstico de diseño y no validación del mundo real.

El paso E0$\rightarrow$E1 requiere instrumentación estable y trazabilidad; E1$\rightarrow$E2, detección reproducible de movimiento y un régimen de diferencias finitas interpretable; E2$\rightarrow$E3, recuperación de forma de onda y evaluación de incertidumbre en datos reservados. Los fallos se documentarán como límites de observabilidad, referencia o modelo.

La validez de esta secuencia depende de que radio, actuador y referencias compartan
una geometría y una escala temporal cuantificadas. El capítulo siguiente traduce
esas condiciones en una arquitectura de medición y un protocolo de sesión.
